% titel: Blatt 4, A2: Energieübergänge im Wasserstoff-Atom
% skript: Mortimer Kap. 6 (Bohr-Modell, Wasserstoffspektrum)

% Teil I: Werkzeuge, kein Lösungsweg.

Energieniveaus und Spektrallinien des Wasserstoffatoms: Jeder Übergang
zwischen zwei Bahnen $n_i \rightarrow n_f$ gehört zu einer diskreten
Energiedifferenz und damit zu einer Spektrallinie.

\begin{werkzeugkasten}[Grundzustand, Übergänge, Spektrallinien]
Grundzustand: kleinstmögliches $n = 1$ (tiefste Energie, stabil).
Bei $n_i > n_f$ wird Energie freigesetzt $\Rightarrow$ \textbf{Emission}
(Licht wird abgegeben); bei $n_i < n_f$ muss Energie aufgenommen werden
$\Rightarrow$ \textbf{Absorption}.

Energiedifferenz und Wellenlänge:
\begin{equation}
  \Delta E = E_{n_f} - E_{n_i}
  = -13.6\left(\frac{1}{n_f^2} - \frac{1}{n_i^2}\right)\ \mathrm{eV},
  \qquad \lambda = \frac{h\,c}{\Delta E}.
\end{equation}
Zahl der Emissionslinien von einem Niveau $n$ aus: alle Paare
$(n_f, n_i)$ mit $n_f < n_i \le n$, also $\binom{n}{2} = n(n-1)/2$.

\textit{Skript:} Mortimer Kap. 6 (Bohr-Modell, Lyman/Balmer/...), Übung 4/2.
\end{werkzeugkasten}

\begin{achtung}[Darauf achten]
\begin{itemize}
  \item Median Vorzeichen: Emission $\to$ $n_f < n_i$; die Energiedifferenz
        wird positiv (freigesetzt) genommen.
  \item Höchste Energie = grösster Abstand der Niveaus (hier $n=1 \leftarrow
        n=4$); niedrigste Energie = kleinster Abstand (hier $n=3 \leftarrow
        n=4$).
  \item Wellenlänge in Metern rechnen (oder eV $\to$ J umwandeln), sonst
        falsche Größenordnung.
  \item Ein Photon niedriger Energie kann einen Grundzustand nur absorbieren,
        wenn es mindestens die Energiedifferenz zum nächsten Niveau
        ($n=1\to n=2$, 10.2 eV) mitbringt.
\end{itemize}
\end{achtung}
